\section{Integral Lewis Distance and Its Parity}
\label{sec:lewis-parity}

The parity of chemical distance depends on the resources represented by the
objective.  A bond-only projection need not have even distance, even for a
balanced closed-shell reaction.  We therefore make the Lewis representation
and its conservation assumptions explicit.

\begin{definition}[Hydrogen-complete closed-shell LLG]
For a reaction side \(X\in\{R,P\}\), let
\(\LLG_X=(V_X,E_X,a_X,b_X)\) be the Lewis-labeled graph of
\cref{sec:occupancy-metrics}, with \(r_X(v)=0\) and integral
\(o_X(u,v)=s_X(u,v)+p_X(u,v)\).  Hydrogens may be explicit vertices or
represented initially by an equivalent hydrogen completion; throughout the
statements below, that completion is understood to be expanded before the
vertex bijection is formed.  ``Closed shell'' means that every nonbonding
electron is included in one of the pairs counted by \(\ell_X\).  Define the
electron-pair inventory
\begin{equation}
  N_{\mathrm{pair}}(\LLG_X)
  =
  \sum_{\{u,v\}\in\binom{V_X}{2}}o_X(u,v)
  +
  \sum_{v\in V_X}\ell_X(v).
  \label{eq:pair-inventory}
\end{equation}
\end{definition}

Let \(\alpha:V_R\to V_P\) preserve the element label \(el\) and any declared
auxiliary isotope tag.  On
closed-shell endpoints, the stored-Kekul\'e electron distance
\cref{eq:kekule-electron-distance} reduces to
\begin{align}
\deK(\alpha)
={}&
\sum_{\{u,v\}\in\binom{V_R}{2}}
\left|
  o_R(u,v)-o_P\bigl(\alpha(u),\alpha(v)\bigr)
\right|
\nonumber\\
&+
\sum_{u\in V_R}
\left|
  \ell_R(u)-\ell_P\bigl(\alpha(u)\bigr)
\right|.
\label{eq:lewis-distance}
\end{align}
Thus a broken bond, a formed bond, a lost lone pair, and a gained lone pair
each contribute one.  Formal charge is not added as a separate term because
it is derived from the same Lewis resources.

Equivalently, let \(B_X\) be the symmetric bond--electron matrix with
\(B_X(u,v)=o_X(u,v)\) for \(u\ne v\) and \(B_X(u,u)=\ell_X(u)\).  If
$\Pi_\alpha$ is the permutation matrix of $\alpha$, then
\begin{equation}
  \deK(\alpha)
  =
  \frac12\sum_{u\ne v}
  \left|
    (B_R)_{uv}
    -
    (\Pi_\alpha^{\mathsf T}B_P\Pi_\alpha)_{uv}
  \right|
  +
  \sum_u
  \left|
    (B_R)_{uu}
    -
    (\Pi_\alpha^{\mathsf T}B_P\Pi_\alpha)_{uu}
  \right|.
  \label{eq:be-distance}
\end{equation}
The factor $1/2$ removes the duplicate entries of each undirected bond; it
does not apply to the diagonal lone-pair loci.

\begin{proposition}[Signed resource identity]
\label{prop:signed-resource}
For an arbitrary integral closed-shell LLG pair
\((\LLG_R,\LLG_P)\) and compatible map \(\alpha\),
let $\delta_k=y_k^\alpha-x_k$ range over all bond and lone-pair loci, and put
\[
  F(\alpha)=\sum_{\delta_k>0}\delta_k,
  \qquad
  C(\alpha)=\sum_{\delta_k<0}(-\delta_k).
\]
Then
\begin{align}
  \deK(\alpha)&=F(\alpha)+C(\alpha),\\
  N_{\mathrm{pair}}(\LLG_P)-N_{\mathrm{pair}}(\LLG_R)
    &=F(\alpha)-C(\alpha).
\end{align}
Consequently,
\begin{equation}
  \deK(\alpha)
  \equiv
  N_{\mathrm{pair}}(\LLG_R)+N_{\mathrm{pair}}(\LLG_P)
  \pmod 2.
  \label{eq:parity-identity}
\end{equation}
In particular, the parity is independent of the AAM.
\end{proposition}

\begin{proof}
The first identity partitions the absolute differences into their positive
and negative parts.  The second follows by summing the signed differences;
transport through $\alpha$ only permutes product resource loci.  Subtraction
and addition are identical modulo two, which gives
\cref{eq:parity-identity}.
\end{proof}

\begin{lemma}[Pair-inventory conservation]
\label{lem:pair-conservation}
Suppose \(\LLG_R\) and \(\LLG_P\) are complete closed-shell Lewis structures,
have the
same element and isotope inventory, and have the same total formal charge.
Then
\(N_{\mathrm{pair}}(\LLG_R)=N_{\mathrm{pair}}(\LLG_P)\).
\end{lemma}

\begin{proof}
Let \(z_X(v)\) be the neutral-atom valence-electron count cached from
\(el_X(v)\), and let \(q_X(v)\) and \(Q_X=\sum_v q_X(v)\) denote local and
total formal charge.
For a complete closed-shell Lewis structure, local electron bookkeeping gives
\[
  q_X(v)
  =
  z_X(v)
  -2\ell_X(v)
  -\sum_{u\ne v}o_X(u,v).
\]
Summing over vertices counts every bond-pair multiplicity twice.  Hence
\begin{equation}
  2N_{\mathrm{pair}}(\LLG_X)
  =
  \sum_{v\in V_X}z_X(v)-Q_X.
  \label{eq:electron-accounting}
\end{equation}
The right-hand side is identical for $R$ and $P$ by atom and charge balance,
and therefore so is \cref{eq:pair-inventory}.
\end{proof}

\begin{theorem}[Evenness of Lewis chemical distance]
\label{thm:lewis-even}
Under the assumptions of \cref{lem:pair-conservation},
\[
  \deK(\alpha)\in2\mathbb N_0
\]
for every compatible atom--atom map $\alpha$.
\end{theorem}

\begin{proof}
Index all bond and lone-pair resource loci in
$\binom{V_R}{2}\uplus V_R$.  Let $x_k$ be the reactant occupancy and let
$y_k^\alpha$ be the product occupancy transported through $\alpha$.  Since
$\alpha$ only permutes product loci,
\[
  \sum_k x_k
  =N_{\mathrm{pair}}(\LLG_R)
  =N_{\mathrm{pair}}(\LLG_P)
  =\sum_k y_k^\alpha.
\]
Set $\delta_k=y_k^\alpha-x_k$.  Then $\sum_k\delta_k=0$, and consequently
\[
  \sum_{\delta_k>0}\delta_k
  =
  \sum_{\delta_k<0}(-\delta_k).
\]
Using \cref{eq:lewis-distance},
\[
  \deK(\alpha)
  =\sum_k|\delta_k|
  =2\sum_{\delta_k>0}\delta_k,
\]
which is an even nonnegative integer.
\end{proof}

Polarity is neither necessary nor sufficient as a mathematical hypothesis.
Every polar reaction represented by hydrogen-complete, balanced, closed-shell
Lewis endpoints with conserved total charge satisfies the theorem.  A polar
reaction with omitted proton/electron inventory, a radical channel, or an
inconsistent endpoint charge does not inherit the claim merely from being
polar.  This distinction is essential when converting mechanism annotations
or RSMI records into LLGs.

Pair-inventory equality is a direct map-independent sufficient condition and
additionally gives $F=C$; closed-shell atom and charge balance in
\cref{lem:pair-conservation} is a chemically transparent sufficient condition
for that equality.  The sharp parity criterion is
\cref{eq:parity-identity}: even unequal endpoint inventories give even
distance when their sum is even, and odd inventory sum gives odd distance
under every map.  An observed odd value under purported theorem conditions is
therefore a diagnostic for an incomplete hydrogen/electron inventory,
fractional bond representation, a radical channel, or an implementation
mismatch; it cannot be repaired by choosing a different AAM.

\begin{corollary}[Exact distance lattice]
\label{cor:distance-lattice}
All compatible AAMs of a fixed balanced closed-shell reaction lie on the
lattice $2\mathbb N_0$.  In particular, if
\(d^\star=\min_\alpha\deK(\alpha)\), then every nonempty distance
shell has the form \(d^\star+2k\) for some \(k\in\mathbb N_0\).
\end{corollary}

\begin{proof}
\Cref{thm:lewis-even} places both \(\deK(\alpha)\) and \(d^\star\) in
\(2\mathbb N_0\).  Their nonnegative difference is therefore \(2k\) for some
\(k\in\mathbb N_0\).
\end{proof}

The conclusion is independent of optimality and symmetry: it holds before
quotienting maps by ITS isomorphism.  It justifies an exact search through
\(\deK(\alpha_0)+2q\) without examining intervening odd thresholds.  It does
not assert that every such shell is nonempty.

\begin{definition}[Bounded ITS neighborhood]
\label{def:bounded-its-neighborhood}
Let \(\mathcal A\) be the set of element-compatible AAMs and let
\(\alpha_0\in\mathcal A\) be the supplied map.  For
\(q\in\mathbb N_0\), define
\begin{equation}
  \mathcal N_q
  =
  \left\{
    [\Upsilon(\LLG_R,\LLG_P,\alpha)]_{\cong}
    \;\middle|\;
    \alpha\in\mathcal A,\quad
    \deK(\alpha)\le
      \deK(\alpha_0)+2q
  \right\}
  \setminus
  \left\{[\Upsilon(\LLG_R,\LLG_P,\alpha_0)]_{\cong}\right\}.
  \label{eq:bounded-its-neighborhood}
\end{equation}
The elements of $\mathcal N_q$, rather than index-distinct AAMs, are the soft
negatives produced by Synister.  Thus two AAMs inducing isomorphic ITS graphs
define the same output.
\end{definition}

\begin{proposition}[Exact cohort form of the local-profile relaxation]
\label{prop:grouped-local-profile}
At a residual state, let $c(r,p)$ be the incident-profile mismatch assigned to
compatible pair $r\mapsto p$, and let the local-profile relaxation be the
minimum of $\sum_r c(r,\sigma(r))$ over residual bijections $\sigma$.  Partition
reactant rows by identical complete domains and incident profiles.  Partition
product columns by identical incident profiles and identical incidence to
those row-group domains.  If $x_{ij}$ transports the row multiplicity of group
$i$ to the column capacity of group $j$, then
\[
  \min_\sigma\sum_r c(r,\sigma(r))
  =
  \min_x\sum_{i,j}c_{ij}x_{ij}
\]
subject to the corresponding integer row and column margins and forbidden
cells.  Therefore dividing the grouped optimum by two gives exactly the same
admissible residual-edge bound as the dense assignment.
\end{proposition}

\begin{proof}
Every dense assignment induces a feasible transport matrix by counting its
row-group--column-group pairs, with identical cost because $c$ is constant on
each allowed cell block.  Conversely, the integral margins of every feasible
transport matrix permit its $x_{ij}$ units to be matched arbitrarily between
the corresponding rows and columns; domain-incidence refinement makes every
such cell allowed.  This expands $x$ to a dense bijection of the same cost.
The two optima are equal.  Each residual unordered edge contributes to the
incident profile of both endpoints, so the standard factor one half remains
admissible and unchanged.
\end{proof}

\begin{proposition}[Domain-compatible residual edge transport]
\label{prop:edge-transport}
At a residual search state, form a bipartite graph whose left vertices are
unmapped reactant edges and whose right vertices are residual product edges.
Join $e_R=\{r,s\}$ to $e_P=\{p,q\}$ exactly when the current domains permit
either orientation of their endpoints, and give the pair weight
$w_b\min\{o_R(e_R),o_P(e_P)\}$.  If $U$ is the maximum weight of a matching
in this graph, then
\[
  w_b\left(\sum_{e_R}o_R(e_R)+\sum_{e_P}o_P(e_P)\right)-2U
\]
is an admissible lower bound on the unmapped bond-distance contribution.
Grouping identical edge rows and columns and solving the resulting integer
transportation problem leaves $U$ unchanged.
\end{proposition}

\begin{proof}
Any residual atom bijection sends distinct reactant edges to distinct product
pairs.  Every image pair that is a product edge therefore defines a matching
between the two residual edge sets, and its preserved bond mass is the sum of
the corresponding pair weights.  It cannot exceed the unrestricted matching
optimum $U$.  The bond $\ell_1$ identity is total reactant mass plus total
product mass minus twice the preserved overlap, which proves the bound.
Identical compatibility neighborhoods and edge orders define constant-cost
complete bipartite blocks; the contraction and integral expansion argument of
\cref{prop:grouped-local-profile} proves equality of the grouped and expanded
matching optima.  Zero-cost dummy margins represent unmatched edges.
\end{proof}

\begin{proposition}[Dual-certified vertex--edge relaxation]
\label{prop:vertex-edge-lp}
At a partial mapping, let $x_{rp}$ be the fractional assignment of residual
reactant atom $r$ to an allowed product atom $p$, with unit row and column
margins.  Let $z$ pair compatible orientations of residual reactant and
product edges, with edge-capacity constraints and endpoint marginals bounded
by the corresponding $x_{rp}$ on both endpoint graphs.  Minimize the exact
linear mapped-boundary cost plus residual endpoint bond mass minus twice the
paired overlap.  The resulting continuous optimum is a lower bound on the
remaining electron distance.

Moreover, let $(y,u)$ be any proposed equality and inequality dual vectors,
with $u\leq0$.  If $v_j$ is the positive violation of dual column $j$, then
\[
 b^Ty+d^Tu-\sum_jv_j
\]
is a certified lower bound when every primal variable lies in $[0,1]$.
Thus evaluating this expression in exact rational arithmetic makes the prune
independent of floating-point solver tolerances.
\end{proposition}

\begin{proof}
Every compatible residual bijection gives an integral feasible point: its
atom variables encode the bijection and its edge variables encode precisely
the product edges reached by mapped reactant edges.  The overlap identity then
makes the objective its exact boundary-plus-interior distance.  Relaxing
integrality can only decrease the minimum.  For the second claim, weak duality
holds after subtracting $v_jx_j$ from every violated dual column; since
$0\leq x_j\leq1$, replacing this by $v_j$ can only decrease the bound.  Exact
evaluation therefore yields a valid repaired-dual certificate.
\end{proof}

\begin{theorem}[Exactness of uncapped production enumeration]
\label{thm:production-exact}
Suppose the branch-and-bound lower bounds are admissible, the hypotheses of
\cref{prop:leaf-twins,prop:lifted-endpoint-symmetry,prop:leaf-transport,%
prop:its-leaf-quotient} hold whenever their respective reductions are invoked,
and no wall-time, node, or class cap is imposed.  Then the production
enumerator terminates and returns exactly one witness for every element of
$\mathcal N_q$.
\end{theorem}

\begin{proof}
Before quotienting, compatibility-class branching spans every bijection in
$\mathcal A$.  A partial branch is removed by a distance bound only when every
completion exceeds the threshold in \cref{eq:bounded-its-neighborhood}.
\Cref{prop:leaf-twins,prop:lifted-endpoint-symmetry,prop:leaf-transport}
replace labeled branches or suffixes by exact orbit representatives and
therefore do not remove an ITS class.
\Cref{prop:dynamic-residual-orbits} replaces candidate children only when an
exact automorphism of the fully colored residual state transports one child
to the other.
Every feasible leaf is scored independently from the complete map.
By \cref{prop:its-leaf-quotient}, the final collision-resolved isomorphism test
merges two leaves if and only if their ITS graphs are isomorphic; the complete
ground-truth class is removed by the same test.  The search space and every
transportation suffix are finite, which proves termination and the claimed
set equality.
\end{proof}

\begin{proposition}[Exact recursive search sharding]
\label{prop:exact-search-sharding}
At any search state, let $r$ be unmapped with complete current product domain
$D(r)$.  If
\[
  D(r)=D_1\mathbin{\dot\cup}\cdots\mathbin{\dot\cup}D_s,
\]
then restricting child $j$ to $\alpha(r)\in D_j$ gives pairwise-disjoint
mapping spaces whose union is the parent mapping space.  Replacing any child
recursively by another such partition preserves this property.  Consequently,
merging the exact ITS quotients of the terminal children returns the same
class set as the unsharded parent if and only if no terminal child is missing
and every terminal child is exhausted.
\end{proposition}

\begin{proof}
Every complete extension assigns $r$ one image $p\in D(r)$.  Disjointness of
the $D_j$ places $p$ in exactly one child, proving both coverage and
uniqueness.  Applying the same argument inductively at an interrupted child
gives a partition tree whose terminal scopes remain disjoint and exhaustive.
The union of their valid mappings is therefore the parent's valid mapping
set.  Exact ITS quotienting commutes with taking this union: it merges
isomorphic witnesses but cannot create or remove a class.  If a terminal
child is missing or incomplete, its mapping set has not been exhausted and
the equality is not certified.
\end{proof}

\begin{proposition}[Exact completed residual-family quotient]
\label{prop:exact-residual-family}
Consider two search states for the same reaction and threshold.  Suppose
(i) their pulled product labels and edges on mapped reactant coordinates are
identical, and (ii) their residual product graphs are isomorphic under a map
that preserves each product resource label, its complete allowed-reactant
domain, every labeled attachment to a mapped reactant coordinate, and every
residual edge label.  Then their valid residual completions are in bijection
and yield the same set of labeled ITS-isomorphism classes and distances.
Consequently, after one such state is exhausted, another state in the same
family may be skipped without losing a unique ITS class.
\end{proposition}

\begin{proof}
Let $\psi$ be the label-preserving residual isomorphism.  If $\beta$ extends
the first state, define the second residual assignment by replacing each
product image $p$ with $\psi(p)$.  Preservation of the complete allowed-row
label makes this assignment feasible and bijective; applying $\psi^{-1}$
gives the inverse construction.  The mapped core is fixed pointwise in
reactant coordinates.  Preserved boundary attachments determine every
mapped--residual pulled edge, and preserved residual labels determine every
remaining node and edge resource.  Corresponding completions therefore have
identical pulled product resource graphs up to relabeling of residual product
identities.  Their signed difference graphs, electron distances, and exact
ITS classes coincide.

The implementation commits a family only after its subtree returns without
truncation.  For forest and unicyclic residual components it compares exact
colored canonical forms obtained from tree centers and dihedral cycle
minimization; more general components use exact labeled graph isomorphism.
WL hashes are not equality witnesses in this reduction.
\end{proof}

\begin{proposition}[Exact dynamic residual candidate orbits]
\label{prop:dynamic-residual-orbits}
At a partial map $\alpha_S$, color every residual product vertex $p$ by its
complete resource label, current allowed-reactant set, and multiset of labeled
attachments to mapped reactant coordinates; retain every residual product
edge and its label.  Let $r$ be the next reactant and let $p,q\in D_S(r)$.
If an automorphism $\psi$ of this colored residual graph satisfies
$\psi(p)=q$, then the children $\alpha_S\cup\{r\mapsto p\}$ and
$\alpha_S\cup\{r\mapsto q\}$ have identical exact ITS-class completion sets.
Thus one candidate per such orbit is sufficient for unique-ITS enumeration.
\end{proposition}

\begin{proof}
For every completion through $p$, compose each remaining product image with
$\psi$.  The allowed-reactant color preserves feasibility and all-different,
while $\psi^{-1}$ gives the inverse map.  The new pulled resource at $r$ is
the same because $p$ and $q$ have equal colors.  Their mapped-boundary
attachments coincide, and $\psi$ transports every edge from the selected
vertex and every edge wholly inside the residual graph.  The pre-existing
mapped core is common to both children.  Corresponding full pulled product
graphs, signed difference graphs, electron distances, and ITS classes are
therefore isomorphic.

For pseudoforest components, equality of the implementation's rooted colored
canonical forms is equivalent to the existence of the required rooted
isomorphism: trees are coded from their centers and unicyclic components are
minimized over both directions and all cycle rotations.  A relevant component
with more than one independent cycle causes a no-pruning fallback; no hash is
accepted as an orbit proof.
\end{proof}

\subsection{Machine-checkable completeness certificates}
\label{sec:certificate-contract}

For one reaction, let $C_i$ be true exactly when endpoint validation succeeds,
the search root is the full element--isotope-compatible bijection domain,
every admissible-bound and symmetry reduction satisfies its stated
hypotheses, every leaf-transport or solver subproblem is exhausted, every
completed residual-family skip satisfies
\cref{prop:exact-residual-family}, every dynamic candidate skip satisfies
\cref{prop:dynamic-residual-orbits}, every
declared shard partition is disjoint and covering, every terminal shard is
present and exhausted,
emitted map and retained witness is independently rescored, WL--3 is used
only as a prefilter before exact labeled isomorphism, and no scope limit
fires.  The implementation sets
\[
  \texttt{class\_complete}_i
  =
  \bigwedge C_i
\]
and signs the normalized case payload.  A timeout may preserve useful
witnesses, but it makes this conjunction false.

For recursive searches, an atomic shard checkpoint may preserve only exhausted
terminal leaves.  Each stored leaf is bound to the reaction, implementation,
search policy, and exact fixed-map/domain scope by digests.  On resume, the
auditor requires the complete and frontier leaves to be a prefix-free cover of
the original root partition before any leaf is reused.  By
\cref{prop:exact-search-sharding}, merging all exhausted leaves after the
frontier becomes empty is exactly the uninterrupted parent search.  An
interrupted frontier leaf contributes neither a completeness claim nor
provisional classes.

\begin{corollary}[Audited cohort completeness]
\label{cor:cohort-certificate}
For a frozen manifest $S$, suppose every stored digest is valid and every
$i\in S$ satisfies $\texttt{class\_complete}_i$.  Then the stored witness set
for every $i$ is exactly its declared non-reference ITS quotient, and the
campaign is class-complete for all $|S|$ scopes.
\end{corollary}

\begin{proof}
Each true case conjunction instantiates the hypotheses and conclusion of
\cref{thm:production-exact} for its declared threshold and label policy.
Manifest equality prevents result-dependent replacement of difficult cases,
while digest validation detects payload corruption.  Taking the conjunction
over $S$ proves the statement case by case.  If even one conjunction is
false, this corollary provides no full-cohort claim.
\end{proof}

\begin{figure}[t]
\centering
\begin{tikzpicture}[x=1cm,y=1cm,font=\sffamily]
  \node[iocard,minimum width=2.35cm,minimum height=1.05cm,align=center]
    (manifest) at (0,0)
    {\textbf{Frozen manifest}\\[-1pt]\scriptsize dataset + scope\\[-2pt]
     \scriptsize hashes};
  \node[iocard,minimum width=2.35cm,minimum height=1.05cm,align=center]
    (root) at (3.15,0)
    {\textbf{Full map root}\\[-1pt]\scriptsize element/isotope\\[-2pt]
     \scriptsize bijections};
  \node[iocard,minimum width=2.35cm,minimum height=1.05cm,align=center]
    (search) at (6.3,0)
    {\textbf{Exact search}\\[-1pt]\scriptsize shards + lower bounds\\[-2pt]
     \scriptsize exact family quotients};
  \node[iocard,minimum width=2.35cm,minimum height=1.05cm,align=center]
    (transport) at (9.45,0)
    {\textbf{Leaf transport}\\[-1pt]\scriptsize DFS or exhaustive\\[-2pt]
     \scriptsize CP--SAT};

  \node[iocard,minimum width=2.35cm,minimum height=1.05cm,align=center]
    (verify) at (9.45,-2.05)
    {\textbf{Independent check}\\[-1pt]\scriptsize map + $\de$\\[-2pt]
     \scriptsize witness};
  \node[iocard,minimum width=2.35cm,minimum height=1.05cm,align=center]
    (wl) at (6.3,-2.05)
    {\textbf{WL--3 bucket}\\[-1pt]\scriptsize necessary invariant\\[-2pt]
     \scriptsize only};
  \node[iocard,minimum width=2.35cm,minimum height=1.05cm,align=center]
    (iso) at (3.15,-2.05)
    {\textbf{Exact quotient}\\[-1pt]\scriptsize labeled ITS\\[-2pt]
     \scriptsize isomorphism};
  \node[iocard,minimum width=2.35cm,minimum height=1.05cm,align=center,
    draw=ioGreen!65!black,fill=ioGreen!5]
    (case) at (0,-2.05)
    {\textbf{Complete case}\\[-1pt]\scriptsize signed witnesses\\[-2pt]
     \scriptsize + empty frontier};
  \node[iocard,minimum width=3.55cm,minimum height=.92cm,align=center,
    draw=ioVermilion!65!black,fill=ioVermilion!4]
    (cohort) at (2.15,-3.7)
    {\textbf{Cohort certificate}\\[-1pt]
     \scriptsize $\bigwedge_{i\in S}\texttt{class\_complete}_i$};
  \node[iocard,minimum width=3.75cm,minimum height=.92cm,align=center,
    draw=ioVermilion!65!black,fill=ioVermilion!4]
    (frontier) at (8.3,-3.7)
    {\textbf{Resumable, incomplete}\\[-1pt]
     \scriptsize signed nonempty frontier};

  \draw[ioflow] (manifest) -- (root);
  \draw[ioflow] (root) -- (search);
  \draw[ioflow] (search) -- (transport);
  \draw[ioflow] (transport) -- (verify);
  \draw[ioflow] (verify) -- (wl);
  \draw[ioflow] (wl) -- (iso);
  \draw[ioflow] (iso) -- (case);
  \draw[ioflow] (case.south) |- (cohort.west);
  \draw[ioflow] (transport.south) -- ++(0,-.35) -| (frontier.north);
  \node[font=\scriptsize,text=black!58,align=center,anchor=south]
    at (8.45,-2.95)
    {timeout / solver unknown / cap};
\end{tikzpicture}

\caption{Proof-carrying unique-ITS campaign in the shared LLG figure style.
Blue arrows are exact transformations; a limit yields a signed, resumable
nonempty frontier, while the final cohort claim is an explicit conjunction of
independently audited empty-frontier cases over the frozen manifest.}
\label{fig:certificate-pipeline}
\end{figure}

\begin{corollary}[Optimal-shell certificate]
\label{cor:optimal-shell-certificate}
Let
\[
  \widehat{\mathcal N}_q
  =
  \mathcal N_q
  \cup\{[\Upsilon(\LLG_R,\LLG_P,\alpha_0)]_{\cong}\}.
\]
After an uncapped run satisfying \cref{thm:production-exact}, the minimum
distance represented in \(\widehat{\mathcal N}_q\) is the global optimum
\(d^\star\).  Moreover, for every
\(d\le \deK(\alpha_0)+2q\), the represented
classes at distance $d$ are exactly all compatible ITS-isomorphism classes at
that distance.
\end{corollary}

\begin{proof}
The supplied map is compatible, hence
\(d^\star\le \deK(\alpha_0)\le \deK(\alpha_0)+2q\).  Thus every global
minimizer satisfies the enumeration bound.  The set equality in
\cref{thm:production-exact}, after restoring the deliberately excluded
ground-truth class, contains every bounded class and no unbounded class.
Restricting that equality to a fixed distance $d$ proves the shell statement.
\end{proof}

Consequently, ``optimal'' and ``complete'' have separate meanings in the
reported data.  The smallest nonempty floor is globally optimal by
\cref{cor:optimal-shell-certificate}; completeness additionally states that
all non-isomorphic alternatives on every floor through the declared bound
have been recovered.  Search limits invalidate the latter certificate and
are therefore never reported as proof-complete output.

\begin{figure}[t]
\centering
\resizebox{\linewidth}{!}{% Lewis-resource parity for ammonia protonation.
% Input-only TikZ file; figures/concept_style.tex is loaded by the parent.
\begin{tikzpicture}[
  font=\sffamily,
  line cap=round,
  line join=round,
  note/.style={font=\scriptsize, text=black!62, align=center},
  good/.style={font=\scriptsize\bfseries, text=ioGreen!48!black},
  minusres/.style={
    resource,
    draw=ioVermilion!55!black,
    fill=ioVermilion!6
  },
  plusres/.style={
    resource,
    draw=ioGreen!45!black,
    fill=ioGreen!6
  }
]

% ---------------------------------------------------------------------------
% A. Lewis structures and endpoint change.
% ---------------------------------------------------------------------------
\begin{scope}[shift={(0,0)}]
  \iosub{-2.30,2.02}{A}{Lewis-complete reaction}
  \begin{pgfonlayer}{background}
    \draw[iopanel] (-2.52,-2.12) rectangle (2.52,2.28);
  \end{pgfonlayer}

  \coordinate (nR) at (-1.58,0.54);
  \coordinate (hR) at (-0.22,0.54);
  \coordinate (clR) at (0.88,0.54);
  \node[catom,text=cdkN] (amineR) at (nR) {H$_3$N};
  \node[celp] at ($(nR)+(0.02,0.34)$) {};
  \node[celp] at ($(nR)+(0.18,0.34)$) {};
  \node[catom,text=cdkH] (protonR) at (hR) {H};
  \node[catom,text=cdkHal] (chlorineR) at (clR) {Cl};
  \node[font=\normalsize] at (-0.88,0.54) {$+$};
  \draw[cbroken] (protonR.east)--(chlorineR.west);
  \node[note,anchor=west] at (1.28,0.54) {$\delta=-1$};

  \coordinate (nP) at (-1.58,-0.72);
  \coordinate (hP) at (-0.22,-0.72);
  \coordinate (clP) at (0.88,-0.72);
  \node[catom,text=cdkN] (amineP) at (nP) {H$_3$N$^{+}$};
  \node[catom,text=cdkH] (protonP) at (hP) {H};
  \node[catom,text=cdkHal] (chlorineP) at (clP) {Cl$^{-}$};
  \node[font=\normalsize] at (0.33,-0.72) {$+$};
  \draw[cformed] (amineP.east)--(protonP.west);
  \node[celp] at ($(clP)+(0.04,0.34)$) {};
  \node[celp] at ($(clP)+(0.20,0.34)$) {};
  \node[note,anchor=west] at (1.28,-0.72) {$\delta=+1$};

  \draw[ioflow] (2.05,0.18) to[bend left=18] (2.05,-0.40);
  \node[note] at (0,-1.55)
    {one bond pair and one lone pair are relocated};
\end{scope}

% ---------------------------------------------------------------------------
% B. Four signed unit resource changes.
% ---------------------------------------------------------------------------
\begin{scope}[shift={(5.35,0)}]
  \iosub{-2.30,2.02}{B}{Signed resource difference}
  \begin{pgfonlayer}{background}
    \draw[iopanel] (-2.52,-2.12) rectangle (2.52,2.28);
  \end{pgfonlayer}

  \node[minusres] (lpN) at (-1.18,0.73)
    {$\operatorname{lp}(N)$\\[-1pt]$-1$};
  \node[plusres] (nh) at (1.18,0.73)
    {$m(N,H)$\\[-1pt]$+1$};
  \node[minusres] (hcl) at (-1.18,-0.73)
    {$m(H,Cl)$\\[-1pt]$-1$};
  \node[plusres] (lpCl) at (1.18,-0.73)
    {$\operatorname{lp}(Cl)$\\[-1pt]$+1$};

  \draw[-{Latex[length=1.8mm]},ioBlue!85!black,line width=.9pt]
    (lpN)--(nh);
  \draw[-{Latex[length=1.8mm]},ioBlue!85!black,line width=.9pt]
    (hcl)--(lpCl);
  \node[note] at (0,-1.58)
    {$\sum\delta^+=2=\sum\delta^-$};
\end{scope}

% ---------------------------------------------------------------------------
% C. Distance and general parity identity.
% ---------------------------------------------------------------------------
\begin{scope}[shift={(10.70,0)}]
  \iosub{-2.30,2.02}{C}{Distance and parity}
  \begin{pgfonlayer}{background}
    \draw[iopanel] (-2.52,-2.12) rectangle (2.52,2.28);
  \end{pgfonlayer}

  \node[iocard,align=center,text width=3.78cm,font=\scriptsize] at (0,0.82)
    {$d_{\mathrm e,K}
      =|{-1}|+|{+1}|$\\[-1pt]
     $\phantom{d_{\mathrm e,K}=}+|{-1}|+|{+1}|$\\[2pt]
     $=4=2\sum_k\delta_k^+$};

  \node[iocard,align=center,text width=3.78cm,font=\scriptsize] at (0,-0.52)
    {$N_{\mathrm{pair}}(\mathfrak L_R)
      =N_{\mathrm{pair}}(\mathfrak L_P)\quad\Longrightarrow$\\[-1pt]
     $d_{\mathrm e,K}=2F\in2\mathbb N_0$};
  \node[iocard,fill=ioGreen!6,draw=ioGreen!30!black,
        align=center,text width=3.78cm,font=\scriptsize\bfseries,
        text=ioGreen!48!black] at (0,-1.45)
    {balanced closed shell\\integral Kekul\'e form};
\end{scope}

\draw[ioflow] (2.65,0)--(2.98,0);
\draw[ioflow] (8.00,0)--(8.33,0);

\begin{scope}[shift={(4.08,-2.46)}]
  \draw[cbroken] (0,0)--(0.50,0);
  \node[anchor=west,font=\scriptsize] at (0.60,0) {decreased};
  \draw[cformed] (2.18,0)--(2.68,0);
  \node[anchor=west,font=\scriptsize] at (2.78,0) {increased};
  \node[celp] at (4.68,-0.03) {};
  \node[celp] at (4.82,-0.03) {};
  \node[anchor=west,font=\scriptsize] at (4.97,0) {lone pair};
\end{scope}
\end{tikzpicture}
}
\caption{Lewis-resource accounting for
\ce{NH3 + HCl -> NH4+ + Cl-}.  The H--Cl bond pair is
removed, the N--H bond pair is formed, the nitrogen lone pair is consumed, and
a chlorine lone pair is gained.  Negative and positive resource changes both
sum to two, hence \(\deK=4\).}
\label{fig:lewis-parity}
\end{figure}

\paragraph{Relation to alternating electron-pair walks.}
Flamm, M\"uller, and Stadler model bonding pairs as multigraph edges and
nonbonding pairs as loops.  For degree-preserving maps, their signed
difference multigraph decomposes into alternating closed walks
\citep{flamm2025electronpushing}.  \Cref{thm:lewis-even} is the corresponding
global resource statement needed here: equality of total pair inventory is
sufficient for parity, whereas alternating closure additionally balances
signed incidences locally at every atom.

\paragraph{Boundary of the theorem.}
Three qualifications are essential.  First, aromatic systems are evaluated
in an integral Kekul\'e realization; the value $1.5$ is not an electron-pair
multiplicity.  Second, omitted hydrogens must be restored or their total
implicit-hydrogen inventory must be balanced.  Third, radical electrons are
not pairs.  Homolytic cleavage has one lost bond pair and two gained radical
electrons, giving \(\deK=2\) but \(\Dres=3\).  Radical reactions therefore
obey a different accounting identity and are excluded from
\cref{thm:lewis-even}; the two objectives must not be conflated.

\subsection{FlowER parity audit and exhaustive spectrum experiment}
\label{sec:lewis-parity-experiment}

We first audited the theorem hypotheses independently of enumeration on the
fixed complexity-stratified FlowER sample.  All 1,000 reactions parsed, and
all had equal total formal charge and implicit-hydrogen count across the two
sides.  Of these, 915 had radical-free endpoints.  Every one of the 915
closed-shell reactions had equal reactant and product pair inventory in
stored Kekul\'e form, and every ground-truth closed-shell distance \(\deK\)
was even.  There were no parity violations.  Among the 85 reactions outside
the theorem because radicals were present, 37 had odd historical
unit-resource distance \(\Dres\); that statistic is not attributed to
\(\deK\).

As a finite exhaustive check of the algebra, we additionally enumerated every
pair of nonnegative integer bond-and-loop resource vectors with equal total
occupancy, together with every atom permutation, for three atoms through
inventory four and four atoms through inventory three.  This covers
1,353,564 distances over 496 resource vectors.  Every distance was even.
This computation tests the implementation over complete finite resource
spaces; the theorem, rather than the finite check, supplies the general
guarantee.

For the enumeration experiment, the production branch-and-bound search uses
the complete objective in \cref{eq:kekule-electron-distance}: the bond terms are
quadratic mapping costs, whereas lone-pair and radical terms are unary
assignment costs.  The frontier linear-assignment bound includes both unary
terms and interactions with fixed atoms; the two interior bounds cover only
unfixed bond interactions and therefore remain admissible.  Stored Kekul\'e
orders determine the objective, while exact ITS isomorphism uses the
aromatic-order endpoint representation.  One witness is retained per
non-ground-truth ITS class.

For the exact comparison, we fix the 898 reactions on which the independent
SLAP graph baseline completed under its prespecified resource limit.  This
cohort definition is made before inspecting Synister's class counts; the 102
baseline timeouts are not silently counted as negative results.  Synister
first enumerated
\[
  \deK(\alpha)\le \deK(\alpha_0)+2
\]
without a wall-time, visited-node, or class-count limit and checkpointed 647
completed cases.  A terminal follow-up used a 600-second per-case wall limit
on the same cohort, reusing 621 prior completions that had already finished
within that budget and evaluating the other 277 cases.  It recorded all 898
outcomes.  A case is included in the exact totals only when its search
exhausted in at least one campaign; bounded truncations and their partial
classes are excluded.  Such a completed case is a proof-complete
neighborhood.  The smaller of the supplied-map distance and all
alternative-witness distances is its certified global minimum: the supplied
map itself is feasible, so no global minimizer can lie above the enumerated
threshold.  Exact leaf-twin
orbit branching and its integer-transport suffix remove factorial
permutations of explicit hydrogens while retaining distinct hydrogen-transfer
patterns, as proved in
\cref{prop:leaf-twins,prop:lifted-endpoint-symmetry,prop:leaf-transport}.
WL--3 is only an isomorphism
prefilter.  Every collision receives an exact isomorphism test on the
equivalence-preserving pendant-twin quotient of
\cref{prop:its-leaf-quotient}.

% Generated by scripts/generate_lewis_exhaustive_tex.py.
\begin{table}[t]
\centering
\small
\caption{Proof-complete Lewis/Kekul\'e enumeration evidence through $D_{\mathrm{GT}}+2$.  Counts combine completed cases from the preserved uncapped checkpoint and the bounded follow-up; truncated cases are excluded.}
\label{tab:lewis-exhaustive}
\begin{tabular}{lrr}
\toprule
Quantity & All completed & Closed-shell completed \\
\midrule
Reactions & 794 & 787 \\
Reactions with alternatives & 791 & 784 \\
Non-GT ITS classes & 347,870 & 346,182 \\
Classes below $D_{\mathrm{GT}}$ & 1,042 & 1,042 \\
Classes at $D_{\mathrm{GT}}$ & 17,041 & 17,020 \\
Classes strictly between GT and GT$+2$ & 0 & 0 \\
Classes at $D_{\mathrm{GT}}+2$ & 329,787 & 328,120 \\
\bottomrule
\end{tabular}
\end{table}

\paragraph{Combined proof-complete evidence.}
\textbf{Preliminary:} the fixed cohort is not fully proof-complete; only cases exhausted in at least one campaign are summarized below. The 600-second follow-up recorded 898 cases: 771 completed and 127 truncated.  The evidence union has 794/898 proof-complete reactions in the fixed baseline-complete FlowER cohort, including 147 added by the follow-up, and 347,870 unique non-ground-truth ITS classes.  Relative to the supplied map, 1,042 classes lie below $D_{\mathrm{GT}}$, 17,041 lie at $D_{\mathrm{GT}}$, and 0 lie strictly between $D_{\mathrm{GT}}$ and $D_{\mathrm{GT}}+2$, while 329,787 lie exactly at $D_{\mathrm{GT}}+2$.  Exactly one representative AAM is stored for each class, so the class total is also the number of generated soft-negative witness records before any downstream sampling.

Across the 787 completed reactions satisfying the theorem hypotheses, 0 enumerated spectra contain an odd distance.  The search visited 34,922,193 states and skipped 2,548,834 leaf-twin branches plus 2,002,029 lifted product-orbit branches.  The 256-automorphism bound was reached in 26 reactions; this remains safe for unique-class completeness but not labeled multiplicities.  The bulk suffix visited 234,583,752 transport states and emitted 849,986 transport-orbit witnesses.  For the 794 cases recorded after activation of the compact WL--3 prefilter, completed-case runtime has median 8.10~s, 90th percentile 155.22~s, and maximum 14755.77~s.

The independent finite resource-space audit evaluated 1,353,564 distances and found 0 odd values.

